\chapter{Loneliness}

\section*{Definition}
Loneliness is a subjective, distressing emotional state arising from a perceived discrepancy between one's desired and actual social relationships. It is distinct from social isolation (objective lack of contact) and can occur even in the presence of others.

\section*{Pathophysiology}
Loneliness activates the hypothalamic--pituitary--adrenal (HPA) axis and sympathetic nervous system, increasing cortisol and inflammatory cytokines (IL-6, CRP). Chronic loneliness impairs prefrontal regulation of the amygdala, heightening hypervigilance to social threat and reducing cognitive reserve.

\section*{Etiology}
\begin{itemize}
  \item \textbf{Psychological factors:} Social anxiety, depression, low self-esteem, negative attribution bias, poor social skills
  \item \textbf{Life transitions:} Bereavement, divorce or separation, retirement, moving to a new city, empty nest, recent immigration
  \item \textbf{Health conditions:} Chronic illness, disability, hearing loss, chronic pain, dementia, terminal illness
  \item \textbf{Social factors:} Living alone, limited social network, caregiving responsibilities, discrimination, poverty
  \item \textbf{Technology/societal:} Excessive social media use, remote work, urbanization, age segregation
\end{itemize}

\section*{Indications for Evaluation}
Seek care if:
\begin{itemize}
  \item Loneliness is persistent, severe, or accompanied by depression, suicidal thoughts, or worsening physical health
  \item There is significant functional decline, social withdrawal, or self-neglect
  \item Loneliness is contributing to unhealthy coping (alcohol misuse, overeating, avoidance of care)
\end{itemize}

\section*{Related Symptoms}
\begin{itemize}
  \item Depression, anxiety, social isolation, grief, low self-esteem, sleep disturbance, fatigue, cognitive decline, substance use
\end{itemize}
\textbf{Note:} Chronic loneliness carries a mortality risk comparable to smoking 15 cigarettes per day and increases cardiovascular and dementia risk.

\section*{Self-Care / Home Management}
\begin{itemize}
  \item Reconnect with existing contacts via phone, video call, or in-person meetings
  \item Join a group based on hobbies (book club, walking group, volunteer organization)
  \item Adopt a pet for companionship and structure (dog for walking, cat for presence)
  \item Limit social media scrolling and focus on meaningful one-on-one interactions
\end{itemize}

\section*{Diagnostic Approach}
\begin{itemize}
  \item Screening using the UCLA Loneliness Scale (Version 3) or De Jong Gierveld Loneliness Scale
  \item Assess for depression (PHQ-9), anxiety (GAD-7), and social isolation (Lubben Social Network Scale)
  \item Evaluate for contributing medical conditions: hearing loss, cognitive impairment, mobility limitations
  \item Consider referral to social work for community resource connection
\end{itemize}

\section*{Treatment / Management Overview}
\begin{itemize}
  \item Cognitive-behavioral therapy (CBT) targeting maladaptive social cognitions and avoidance behaviors
  \item Social prescribing: linking to community groups, volunteer opportunities, and peer support
  \item Treat underlying depression or anxiety with therapy and/or SSRIs
  \item Address sensory deficits (hearing aids, glasses) and mobility barriers (transportation assistance)
  \item Befriending services and telephone support lines for housebound individuals
\end{itemize}

\clearpage
